% ID: SEL-D0069-OB-E02
% ESTADO: cotejada. Segunda pasada visual y borrador compilado revisados.
% Fuente: Selectividad/sel_2009_matematicas_aplicadas/mat apli CC SS03.pdf
% SHA-256: 8cf5f487f300c573430722e40f741601da1756c5e8c13e541f697d1620373545
% Página: 2; opción B; ejercicio original 2.

Un estudio acerca de la presencia de gases contaminantes en la atmósfera de una ciudad indica que el nivel de contaminación viene dado por la función:
\[
C(t)=-0.2t^2+4t+25,\quad0\leq t\leq25\quad(t=\text{años transcurridos desde el año 2000}).
\]
\begin{enumerate}[label=\alph*)]
\item \textbf{(1 punto)} \textquestiondown{}En qué año se alcanzará un máximo en el nivel de contaminación?
\item \textbf{(1 punto)} \textquestiondown{}En qué año se alcanzará el nivel de contaminación cero?
\item \textbf{(1 punto)} Calcule la pendiente de la recta tangente a la gráfica de la función $C(t)$ en $t=8$. Interprete el resultado anterior relacionándolo con el crecimiento o decrecimiento.
\end{enumerate}
